\documentclass[11pt]{article}
\newcommand{\LabNumber}{4}
\newcommand{\LabTitle}{Doubly Linked Lists}
\newcommand{\LabTopic}{Sentinel nodes, bidirectional linking, addFirst/addLast/remove}
\input{common.tex}

\begin{document}
\labheader

\section*{Objectives}
\begin{itemize}[leftmargin=*,nosep]
  \item Understand why sentinel (header/trailer) nodes simplify boundary conditions.
  \item Implement a generic \texttt{DoublyLinkedList<E>} with \texttt{prev} and \texttt{next} pointers.
  \item Implement \texttt{addFirst}, \texttt{addLast}, and \texttt{remove(Node)} in $O(1)$ time.
  \item Compare the trade-offs of singly vs.\ doubly linked lists.
\end{itemize}

\begin{summary}[Quick Reference]
\textbf{Node class.}
\begin{lstlisting}
private static class Node<E> {
    E element;
    Node<E> prev, next;
    Node(E e, Node<E> p, Node<E> n) { element = e; prev = p; next = n; }
}
\end{lstlisting}

\textbf{Sentinel setup.}
\begin{lstlisting}
header  = new Node<>(null, null, null);
trailer = new Node<>(null, header, null);
header.next = trailer;
\end{lstlisting}

\textbf{addBetween} (private helper):
\begin{lstlisting}
Node<E> node = new Node<>(e, predecessor, successor);
predecessor.next = node;
successor.prev  = node;
size++;
\end{lstlisting}

\textbf{remove(Node node):}
\begin{lstlisting}
Node<E> p = node.prev, s = node.next;
p.next = s;  s.prev = p;
size--;
return node.element;
\end{lstlisting}
\end{summary}

\subsection*{Quick Experiments (Try It)}

\begin{tryit}{Sentinels eliminate special cases}
Add a print inside \texttt{addBetween} to confirm the same code path handles both the
empty and non-empty cases. How many \texttt{if}-statements does \texttt{addBetween} need?
\end{tryit}

\begin{tryit}{Backward traversal}
Traverse your list from \texttt{trailer.prev} back to \texttt{header.next}.
Verify the output is the reverse of the forward traversal.
\end{tryit}

\section*{In-Lab Exercises (target: 2 hours)}

\begin{labexercise}{Implement \texttt{DoublyLinkedList<E>} \hfill {\small \itshape (50 min)}}
Build the class using header/trailer sentinels:
\begin{itemize}[leftmargin=*,nosep]
  \item \texttt{int size()}, \texttt{boolean isEmpty()}.
  \item \texttt{E first()}, \texttt{E last()} --- null if empty.
  \item \texttt{void addFirst(E e)}, \texttt{void addLast(E e)}.
  \item \texttt{E removeFirst()}, \texttt{E removeLast()} --- both $O(1)$.
  \item \texttt{String toString()} --- e.g.\ \texttt{[A <-> B <-> C]}.
\end{itemize}
Use private \texttt{addBetween} and \texttt{remove(Node)} helpers.
Test with at least 10 operations including removes from both ends.
\end{labexercise}

\begin{labexercise}{Add at Position \hfill {\small \itshape (25 min)}}
Add \texttt{void add(int index, E e)} that inserts at zero-based position \texttt{index}.
Traverse to node at position \texttt{index}, then call \texttt{addBetween}.
What is the time complexity? When is this worse than an array?
\end{labexercise}

\begin{labexercise}{Deque Adapter \hfill {\small \itshape (25 min)}}
Wrap \texttt{DoublyLinkedList<E>} in a \texttt{LinkedDeque<E>} with:
\texttt{addFirst}, \texttt{addLast}, \texttt{removeFirst}, \texttt{removeLast},
\texttt{peekFirst}, \texttt{peekLast}, \texttt{size}, \texttt{isEmpty} --- all $O(1)$.
\end{labexercise}

\begin{labexercise}{Palindrome Checker \hfill {\small \itshape (20 min)}}
Using \texttt{LinkedDeque<Character>}, write \texttt{boolean isPalindrome(String s)}:
add each character to the deque, then repeatedly remove from both ends and compare.
Test with \texttt{"racecar"}, \texttt{"hello"}, and \texttt{"madam"}.
\end{labexercise}

\begin{aicollab}
\textbf{Suggested prompts:}
\begin{itemize}[leftmargin=*,nosep]
  \item ``Explain how sentinel nodes eliminate edge-case checks in a doubly linked list.''
  \item ``Why is \texttt{removeLast} $O(1)$ in a doubly linked list but $O(n)$ in a singly linked list?''
\end{itemize}
\medskip\textbf{Verify:} Ask the AI to draw pointer changes for \texttt{addFirst} and \texttt{removeLast}
on a one-real-node list. Run those operations and confirm the sentinel pointers match.
\end{aicollab}

\takehomebanner

\begin{trace}[Pointer trace with sentinels]
Draw the doubly linked list (including sentinels H and T) after each step.
Show all \texttt{prev}/\texttt{next} arrows.
\begin{lstlisting}
DoublyLinkedList<String> L = new DoublyLinkedList<>();
L.addFirst("B"); L.addFirst("A"); L.addLast("C"); L.removeFirst(); L.removeLast();
\end{lstlisting}
\end{trace}

\begin{writecode}[Reverse in-place]
Write \texttt{void reverse()} for \texttt{DoublyLinkedList}: swap \texttt{prev} and \texttt{next}
for every node including sentinels. No new nodes created.
\end{writecode}

\begin{drawex}{Singly vs.\ doubly linked list}
Draw side-by-side memory diagrams of a three-element singly linked list and a
three-element doubly linked list (with sentinels). Label every pointer field.
Circle the extra fields in the doubly linked version and state the per-node memory overhead.
\end{drawex}

\vspace{4pt}
\noindent\textbf{Submission.} Save files as \texttt{lab4\_ex*.java} and submit on the course site.

\end{document}
